\chapter{The working habits, and learning from history}
\label{ch:habits}

The default loop ships with habits: small Scheme layers, each written
after a failure seen in real logs, measured on a coding benchmark and kept
because they paid their way (\adr{32}). Reading them is the best
advanced Scheme course in the repository, and knowing them tells you what
your own strategy is wrapping.

\section{The working rules}

\code{working-rules} (in \file{prompts.scm}) is added to both system
prompts. It is the text form of the habits: run command lines through
\code{bash -c}; start a program that keeps running or may hang with
\code{shell-start} instead of waiting on \code{shell-exec}; find the project's own interpreter and environment before
running its code, and never patch the standard library to make tests run;
read with the file tools, not with scripts; ask for the calls that do not
depend on each other in one reply, since every reply sends the whole
conversation again; treat a plan written before
anything ran as a first guess; keep \textbf{one check script} in
\code{\$TMPDIR} that shows the problem before the fix and is rerun
unchanged after every change; follow the conclusion of a request that
weighs options; widen a fix in the general form, not just for the example;
run whole test modules, and check code a browser runs with the project's
browser checks, since a build or a syntax check only shows it is well
formed; compare failures against the unchanged code on a copy of the last
commit (\code{git archive HEAD} into \code{\$TMPDIR/base}), never with
\code{git stash}, and only when tests fail; bisect regressions file by file;
and keep what a later session would need and could not quickly find in the
project notes, with \code{project-notes-save}. In a git work tree,
\code{(git-rule)} adds one more: look at the change with \code{git-status}
and \code{git-diff} before finishing, and leave nothing in it that nothing
uses. In a task started in a git worktree of its own,
\code{(worktree-section)} names its branch, says other tasks may work in
other worktrees meanwhile, and says again that \code{git stash} is not
used there, since every worktree shares one stash (\adr{34}).

If your project has its own rules, add them after these rather than
replacing them (Chapter~\ref{ch:checks}).

\section{The habit files}

Each file wraps \code{run-calls} (and sometimes \code{fold-result},
\code{sub-agent}, \code{model-turn} or \code{nodwyn-turn}) with the idiom from
Chapter~\ref{ch:strategies}, so your \code{init.scm} can wrap them again,
tune their settings, or redefine them away.

\subsection{\file{shell.scm}: command lines, and results a model can read}

A model used to a bash tool writes \code{cd src \&\& pytest | tail -1}.
Through \code{shell-exec}'s program-plus-arguments interface that is
nonsense, and on some tasks a model sent 8 to 15 such malformed calls. So
a call whose arguments contain shell operators is answered, unrun, with
the exact \code{bash -c} call that would run it as meant. A strategy may
skip a model's call but never rewrite it, so this is the honest fix: tell
the model the right call.

Results are shown as plain text (the exit status, then the output), and
long output keeps its start and its end, where test runners print their
summary: \code{shell-head-chars} (1,500) and \code{shell-tail-chars}
(6,000). And \code{ran-ok?} counts a test run with failures as a run.

\subsection{\file{hygiene.scm}: fold what stopped mattering}

About 70\% of the tool output in the main context was test runs and file
reads, and 43\% of the bytes read were from files never edited. Three
rules, main context only:

\begin{enumerate}
  \item a test command that never ran a test is folded to one line once a
    later one does. A command ran tests when it exited 0 or its output
    says so in a runner's words (\code{hygiene-ran-words}: pytest's, Go's
    \code{ok} and \code{-{}-{}- FAIL}, mocha's, jest's, cargo's, and a
    repository validate's summary); a Go package that did not build ran
    none;
  \item a read of a file edited since is \emph{stale}: after
    \code{hygiene-stale-age} (2) rounds it folds to its outline, noting the
    file changed;
  \item a read \code{hygiene-age} (8) rounds old folds to its outline.
\end{enumerate}

Neither of the last two folds a read of a file edited in the last
\code{hygiene-editing-rounds} (4) rounds: the file is being worked on, and
\code{fs-edit} answers with the lines it changed, so its reads stay.

An outline keeps the definition lines (\code{def}, \code{class},
\code{func}, \code{fn}\ldots) with their line numbers, so the model can
read any part back with a ranged \code{fs-read}. Reads fold in batches of
at least \code{hygiene-min-bytes} (24,000), because every fold makes the
server re-read the prompt from that point. With a provider whose prompt
cache works, that re-read costs the full input price for everything after
the first folded read, where a cached token costs a fraction of it, so a
batch also has to pay for itself. It folds when it is at least as large as
the rest of the window after it (\code{(context-tokens result)} measures
that), when the provider's cache is not serving the context anyway (it read
less than \code{hygiene-cache-share}, a quarter, of the recent prompts), or
when the context is past \code{hygiene-pressure} (half) of the high water,
where room matters more than the cache. Otherwise the reads wait. Each
fold's log note says which reason applied.

\subsection{\file{readdef.scm}: \texttt{read-definition}}

Fifteen of sixteen Python scripts that opened a source file were finding a
definition with a regular expression and printing it. \code{read-definition}
(path, name) is that, in one call: the whole function, method or class and
its line range, or for a long class its start and members. And a script
that only reads a file to print it gets a hint naming the file tools, at
most \code{readdef-hints} (3) times a session.

\subsection{\file{envhint.scm}: find the project's environment}

A command that fails with \code{No module named \ldots} gets a hint
appended to its result (with \code{fold-result}): the project most likely
has its own environment, so look for \file{.venv}, conda, tox or nox, and
use its interpreter. At most \code{envhint-hints} (2) times a session. It is
a nice example to read first: twenty lines, one wrapper, one counter.

\subsection{\file{checkloop.scm}: \texttt{second-look}}

When the model's check script does not reproduce a reported problem, the
tool \code{second-look} gets two read-only helpers. The first reads the
request \emph{alone}, literally, and works out which inputs would produce
the outputs it shows. The second compares that reading with the check
script and lists concrete differences to try. The order matters: a helper
shown the check first adopted the check's (wrong) reading. A history search
(\code{git log}, \code{git blame}) before any code edit gets a hint to use
it instead.

\begin{tip}
The comments at the top of each habit file cite the log evidence behind
it, down to task names. When you tune a habit for your project, read those
comments first: they tell you what you would be trading away.
\end{tip}

\section{Tuning and removing}

\begin{scm}
;; Keep reads around longer in a project with big, stable files.
(define hygiene-age 16)

;; Show more of the start of long outputs.
(define shell-head-chars 4000)

;; A project whose tests run with ./check.
(define test-command-words (cons "./check" test-command-words))
\end{scm}

Most habits have a setting that switches them off cleanly:

\begin{scm}
(define envhint-hints 0)        ; no environment hints
(define readdef-hints 0)        ; no hints about scripts that read files
(define second-look-hints 0)    ; no hints to use second-look
(define (hygiene-sweep) #f)     ; never fold old reads
\end{scm}

That last line works because \file{hygiene.scm}'s wrapper calls
\code{hygiene-sweep} by name on every round, and a global name is looked up
when it is called.

\begin{careful}
You can also restore what a habit wrapped, but only for the \emph{last}
layer: \code{(define run-calls checkloop-base-run-calls)} removes
\file{checkloop.scm}'s wrapper, the outermost one. Restoring an inner
layer's base (say \code{envhint-base-run-calls}) would drop every layer
loaded after it as well. Prefer the settings above.
\end{careful}

\section{Learning from history}

A strategy can read this folder's past sessions:

\begin{scm}
(session-list 2)
;=> (((id . "20260930T200053Z-d5a648e8") (first-message . "") (current . #t) ...)
;    ((id . "20260930T200043Z-e0550e21") (first-message . "") (current . #f) ...))

(session-events "20260930T200043Z-e0550e21" '("user.msg" "tool.result"))
;=> each event: ((seq . 2) (ts . "...") (kind . "user.msg") (payload . (...)))
\end{scm}

Only this folder's sessions are listed, and the events come back redacted
and tagged \code{external:log}, like any tool output.

\subsection{\texttt{/optimize}}

The optimizer is user code over these primitives, in
\file{optimize.scm}. Send a message starting with \code{/optimize}:

\begin{txt}
/optimize                 every earlier session of this folder
/optimize 5               the latest five
/optimize ID ID ...       the sessions named
/optimize 5 test runs     any other words are a focus
\end{txt}

It digests the sessions (requests, answers, calls, failures and refusals,
judges' verdicts, sub-agents, notes) with code much like the word counter
in Chapter~\ref{ch:playground}, and hands the digest to a sub-agent
labelled \code{optimizer}, which may read the folder and write only under
\file{.nodwyn/}. Every write asks you, with the diff (a protected path),
so \code{/optimize} needs a person at the prompt. When it finishes,
\code{init.scm} is kept as a proposal until you review its new revision with
\code{nodwyn trust}, then reload or resume. The new code has not run. For
an explicit Go-host init, reload still applies immediately; other reload failures
restore the old file.

\begin{tip}
Run \code{/optimize} after a week of real work in a repository, with a
focus: \code{/optimize 10 how tests are run}. Review its diff like a
colleague's pull request. What it writes is ordinary Scheme you now know
how to read.
\end{tip}

\begin{puzzle}[ 18: your own habit]
Write a habit that appends a hint to the result of any \code{shell-exec}
whose output contains \code{Permission denied}, suggesting the model check
whether the file is executable, at most twice per session. Model it on
\file{envhint.scm}.\ans{18}
\end{puzzle}
